% Конспект и немного воображения
\ticket{04}{Многомерная формула Тейлора (с остатком в форме Лагранжа и Пеано)}
\label{ticket:04}

\begin{theorem}[Формула Тейлора]
Пусть $E\subset\mathbb R^m$ открыто, $f\colon E\to\mathbb R$,
$f\in C^{r+1}(E)$, $x\in B(a,R)\subset E$.
Тогда существует $t\in(0,1)$ такое, что
\[
  \begin{aligned}
    f(x)={}&\sum_{|j|\le r}\frac{1}{j!}
      \frac{\partial^{|j|}f}{\partial x_j}(a)(x-a)^j\\
    &+\sum_{|j|=r+1}\frac{1}{j!}
      \frac{\partial^{r+1}f}{\partial x_j}\bigl(a+t(x-a)\bigr)(x-a)^j.
  \end{aligned}
\]
\end{theorem}

\begin{proof}
Положим
\[
  x=a+h,\qquad h=x-a,\qquad \varphi(t)=f(a+th).
\]
По одномерной формуле Тейлора существует $t\in(0,1)$ такое, что
\[
  \begin{aligned}
    \varphi(1)={}&\varphi(0)+\frac{\varphi'(0)}{1!}1^1+\ldots
      +\frac{\varphi^{(r)}(0)}{r!}1^r\\
    &+\frac{\varphi^{(r+1)}(t)}{(r+1)!}1^{r+1}.
  \end{aligned}
\]
Подставляя производные $\varphi$ по
\hyperref[ticket:03]{лемме о дифференцировании «сдвига»}, получаем требуемую формулу.
\end{proof}

\textbf{Остаток в форме Пеано.}
При тех же условиях, оценивая остаток Лагранжа, получаем
\[
  f(a+h)=\sum_{|j|\le r}\frac{1}{j!}
    \frac{\partial^{|j|}f}{\partial x_j}(a)h^j
    +o(|h|^r),\qquad h\to0.
\]

Немного подробнее: производные порядка $r+1$ непрерывны,
поэтому ограничены в некоторой окрестности $a$.
Следовательно, при $h\to0$
\[
  \sum_{|j|=r+1}\frac{1}{j!}
    \frac{\partial^{r+1}f}{\partial x_j}(a+th)h^j
  =O(|h|^{r+1})=o(|h|^r).
\]
